\scanpage{129}
\begin{center}- 135 -\end{center}
\margin{n\textsuperscript{o} 316}

a) L'algèbre $\widetilde U$ et l'application $\widetilde{\unsure{\varrho|\rho}}$ sont solutions d'un problème d'application universelle: pour tout p-homomorphisme $f$ de $\unsure{g|\mathfrak g}$ dans une algèbre associative $B$ sur $K$ (considérée comme p-algèbre de Lie), il existe un K-homomorphisme et un seul $f'$ de $\widetilde U$ dans $B$ (pour les structures d'algèbres associatives) tel que $f=f'\circ\widetilde{\unsure{\varrho|\rho}}$.

b) Montrer que si $(b_\lambda)_{\lambda\in\Lambda}$ est une base de $\unsure{g|\mathfrak g}$ ($\Lambda$ étant totalement ordonné), $\widetilde{\unsure{\varrho|\rho}}$ est injective, et que, si on identifie $x$ et $\widetilde{\unsure{\varrho|\rho}}(x)$ pour $x\in\unsure{g|\mathfrak g}$, les éléments $\prod_\lambda b_\lambda^{\nu_\lambda}$ (où les $\lambda$ vont en croissant, où les $\nu_\lambda$ sont nuls sauf un nombre fini d'entre eux, et où $0\leq\nu_\lambda<p$ pour tout $\lambda$) forment une base de $\widetilde U$. (Identifiant canoniquement $\unsure{g|\mathfrak g}$ à un sous-module de $U$, on pose $\varphi(x)=x^p-x^{[p]}$ pour tout $x\in\unsure{g|\mathfrak g}$. Pour tout indice composé $\alpha=(\alpha_\lambda)\in\mathbf N^{(L)}$, soit $\alpha_\lambda=\beta_\lambda+p\gamma_\lambda$ avec $0\leq\beta_\lambda<p$, et soit $T_\alpha=(\prod_\lambda b_\lambda^{\beta_\lambda})(\prod_\lambda(\varphi(b_\lambda))^{\gamma_\lambda})$. Montrer que les $T_\alpha$ forment une base de $U$ et les $T_\alpha$ tels que $\gamma=(\gamma_\lambda)\ne0$ une base de $J$; on observera que les $\varphi(b_\lambda)$ appartiennent au centre de $U$).

7) - Soit $\unsure{g|\mathfrak g}$ une p-algèbre de Lie sur un anneau $K$ de caractéristique $p>0$ ($p$ premier). On dit qu'une dérivation $D$ de $\unsure{g|\mathfrak g}$ est une \uline{p-dérivation} si on a $D(x^p)=(\mathrm{ad}\,x)^{p-1}.dx$ pour tout $x\in\unsure{g|\mathfrak g}$. Toute dérivation intérieure est une p-dérivation.

a) Si $L$ est une K-algèbre associative, toute dérivation de $L$ est une p-dérivation lorsque $L$ est considérée comme p-algèbre de Lie. (Utiliser la formule (2) de l'exerc. 19 du § 1).

b) Supposons que $\unsure{g|\mathfrak g}$ ait une base. Pour qu'une dérivation de $\unsure{g|\mathfrak g}$ soit une p-dérivation, il faut et il suffit qu'elle puisse se prolonger en une dérivation de l'algèbre enveloppante restreinte de $\unsure{g|\mathfrak g}$.
\margin{128}
